\section{Diagrama de clases}

El diagrama de clases UML sirve para mostrar una vista estática de las clases de un sistema, sus atributos, métodos y relaciones entre objetos. Sirve para proporcionar una información general sobre la aplicación \cite{clase}. 

El diagrama de clases se compone de tres partes: 
\begin{itemize}[nosep, noitemsep]
	\item \textbf{CUDA}: La figura \ref{fig:class_cuda} muestra la definición de las clases del autómata celular y el autómata celular neuronal. 
	\item \textbf{Modelos}: Son los modelos construidos con pybind11 los cuales son manipulados por ``SimulationController'', como se muestra en la figura \ref{fig:class_models}.
	\item \textbf{UI}: Elementos y controladores de la interfaz gráfica en PyQt5. Ver figura \ref{fig:class_ui}.
\end{itemize}

\begin{figure}[H]
	\centering
	\includegraphics[width=.7 \linewidth]{"./img/classDiagram_cuda.png"}
	\caption{Diagrama de clases: Capa CUDA}
	\label{fig:class_cuda}
\end{figure}

\begin{figure}[H]
	\centering
	\includegraphics[width=.95 \linewidth]{"./img/classDiagram_models.png"}
	\caption{Diagrama de clases: Capa de modelos}
	\label{fig:class_models}
\end{figure}

\begin{figure}[H]
	\centering
	\includegraphics[width=.95 \linewidth]{"./img/classDiagram_ui.png"}
	\caption{Diagrama de clases: Capa UI}
	\label{fig:class_ui}
\end{figure}